\begin{proof}
Put\footnote{Editors' note: In what follows, the original argument has
been expanded, and equality \((1)\) has been corrected.}
\[
 M(w)
 =\operatorname{Card}\bigl(R^+\cap w^{-1}(-R^+)\bigr)
 =\operatorname{Card}\{\gamma\in R^+\mid w(\gamma)\in-R^+\}.
\]
If \(M(w)=0\), then \(w(R^+)=R^+\), so
\(w=1\) by Corollary~\ref{sga3:XXI:5.4}; the theorem is then
trivial: \(m=0\), and assertions (iii)--(v) are empty.

We argue by induction on \(M(w)\). Suppose \(M(w)>0\). There exists
\(\beta_0\in\Delta\) such that \(w(\beta_0)\in-R^+\). Put
\(\alpha_0=\alpha\), and consider
\[
 A=w^{-1}(R^+)\cap R_{\alpha_0,\beta_0}.
\]
This is a system of positive roots of
\(R_{\alpha_0,\beta_0}\). Hence there exists
\(w_0\in W_{\alpha_0,\beta_0}\) such that
\[
 w_0^{-1}\bigl(R^+_{\alpha_0,\beta_0}\bigr)=A.
\]
Put \(w'=ww_0^{-1}\). Lemma~\ref{sga3:XXI:3.4.9} gives
\[
 R^+-R^+_{\alpha_0,\beta_0}
 =
 w_0\bigl(R^+-R^+_{\alpha_0,\beta_0}\bigr),
\]
and therefore
\begin{equation}
\tag{1}
\label{sga3:XXI:5.6:outside-rank-two}
\begin{split}
 &(R^+-R^+_{\alpha_0,\beta_0})\cap w^{-1}(-R^+)\\
 &\qquad=
 w_0\Bigl((R^+-R^+_{\alpha_0,\beta_0})
             \cap w'^{-1}(-R^+)\Bigr).
\end{split}
\end{equation}
On the other hand,
\begin{equation}
\tag{2}
\label{sga3:XXI:5.6:beta-negative}
 \beta_0\in
 R^+_{\alpha_0,\beta_0}\cap w^{-1}(-R^+).
\end{equation}
Since
\(w_0(R_{\alpha_0,\beta_0})=R_{\alpha_0,\beta_0}\), one has
\[
 w_0(-A)
 =R_{\alpha_0,\beta_0}\cap w'^{-1}(-R^+),
\]
and hence
\begin{equation}
\tag{\(2'\)}
\label{sga3:XXI:5.6:rank-two-empty}
\begin{split}
 R^+_{\alpha_0,\beta_0}\cap w'^{-1}(-R^+)
 &=
 R^+_{\alpha_0,\beta_0}\cap w_0(-A)\\
 &=
 R^+_{\alpha_0,\beta_0}\cap(-R^+_{\alpha_0,\beta_0})
 =\varnothing.
\end{split}
\end{equation}
It follows from \eqref{sga3:XXI:5.6:outside-rank-two},
\eqref{sga3:XXI:5.6:beta-negative}, and
\eqref{sga3:XXI:5.6:rank-two-empty} that
\[
 M(w')<M(w).
\]

Put \(\alpha_1=w_0(\alpha_0)\). We show that
\(\alpha_1\in\Delta\), or equivalently that
\(\alpha_0\in w_0^{-1}(\Delta)\). Since
\(w(\alpha_0)\in\Delta\), we have
\[
 \alpha_0\in
 w^{-1}(\Delta)\cap R_{\alpha_0,\beta_0}.
\]
Thus \(\alpha_0\) is a simple root of
\[
 A
 =w^{-1}(R^+)\cap R_{\alpha_0,\beta_0}
 =w_0^{-1}(R^+_{\alpha_0,\beta_0}),
\]
and consequently belongs to
\[
 w_0^{-1}\bigl(\Delta\cap R^+_{\alpha_0,\beta_0}\bigr)
 =
 w_0^{-1}\{\alpha_0,\beta_0\}
\]
by Lemma~\ref{sga3:XXI:3.4.8}. Hence
\(\alpha_1=w_0(\alpha_0)\) equals \(\alpha_0\) or \(\beta_0\).
If \(\alpha_1\ne\alpha_0\), then
\(\alpha_1=\beta_0\) and
\(w_0\in W_{\alpha_0,\alpha_1}\).

Finally,
\[
 \beta=w'(\alpha_1)
 \qquad\text{and}\qquad
 M(w')<M(w),
\]
so the induction hypothesis completes the proof.
\end{proof}
